% =========================================================
% Proof: Interior-Boundary-Exterior Decomposition
% Source: notes/open-closed-sets/open-closed-sets.tex
% Difficulty: Easy
% =========================================================

\subsection*{Interior--Boundary--Exterior Decomposition}
\label{prf:ibe-decomp}

\begin{remark*}[Return]
$\leftarrow$ \hyperref[thm:ibe-decomp]{Back to Theorem in Notes}
\end{remark*}

\begin{proof}
\Claim Let $(X,d)$ be a metric space and $A \subseteq X$.
The interior $\operatorname{int}(A)$, boundary $\partial A$, and exterior
$\operatorname{ext}(A)$ are pairwise disjoint and their union is $X$.

\Given A metric space $(X,d)$ and $A \subseteq X$.
Recall: $x \in \operatorname{int}(A)$ iff some ball around $x$ is $\subseteq A$;
$x \in \operatorname{ext}(A)$ iff some ball around $x$ is $\subseteq X \setminus A$;
$x \in \partial A$ iff every ball around $x$ meets both $A$ and $X \setminus A$.

\Goal Show: (1) the three sets are pairwise disjoint; (2) every $x \in X$ is in exactly one.

\Strategy For any $x \in X$, any ball $B_r(x)$ either
(a) lies entirely in $A$, (b) lies entirely in $X \setminus A$, or (c) meets both.
At least one of these must hold; no two can hold simultaneously.

% -------------------------------------------------------
% YOUR PROOF HERE
% -------------------------------------------------------

\end{proof}

\begin{remark*}[Proof shape]
Exhaustive case analysis on the behaviour of balls around $x$.
Each of the three cases is mutually exclusive, and one must always occur.
\end{remark*}

\begin{remark*}[Dependencies]
\begin{itemize}
  \item Definition~\ref{def:interior}: $\operatorname{int}(A)$.
  \item Definition~\ref{def:exterior}: $\operatorname{ext}(A)$.
  \item Definition~\ref{def:boundary}: $\partial A$.
\end{itemize}
\end{remark*}
